\slcol{अर्जुन उवाच —\\
\Index{मदनुग्रहाय परमं} गुह्यमध्यात्मसंज्ञितम् । \\
यत्त्वयोक्तं वचस्तेन मोहोऽयं विगतो मम ॥ ११.१ ॥ }
\translate{Arjuna uvāca — \\
madanugrahāya paramaṁ guhyamadhyātmasaṁjñitam | \\
yattvayōktaṁ vacastēna mōhō:'yaṁ vigatō mama || 11.1 || }
\cquote{Arjuna said: Out of compassion for me, You have spoken the supreme secret known as \emph{adhyātma}, knowledge of Self. By that word of Yours, my delusion has been dispelled.}
\slcol{\Index{भवाप्ययौ हि भूतानां} श्रुतौ विस्तरशो मया । \\
त्वत्तः कमलपत्राक्ष माहात्म्यमपि चाव्ययम् ॥ ११.२ ॥ }
\translate{bhavāpyayau hi bhūtānāṁ śrutau vistaraśō mayā | \\
tvattaḥ kamalapatrākṣa māhātmyamapi cāvyayam || 11.2 || }
\cquote{O Lotus-eyed (Krishna), I have indeed heard from You in detail of the origin and dissolution of beings, and also of Your inexhaustible greatness.}

\newpage
\begin{mananam}{\mananamfont {Mananam Śloka - 1}}
\normalfont\mananamtext Have there been moments when intense emotions, such as anger, hatred, fear, or panic, clouded my judgment? In such situations, how did gaining the right understanding help me manage or overcome these emotions?
What was my attitude toward the person or object that triggered feelings of hatred or fear before I achieved understanding, and then, after gaining understanding?
What or who was the source of this new understanding that brought about such a transformative effect? How can I access or connect with the source of this understanding in order to experience its benefits consistently?
\end{mananam}
\sreflect
\begin{inspiration}{\inspirationTitleFont{Inspiration}}
\normalfont\mananamtext The words of great beings have a profound ability to uplift and inspire us. At times, we may find ourselves internally blinded by misguided perspectives and tendencies, unable to perceive life in its true essence. However, with the awakening of understanding, we come to realize that we have been dwelling in darkness, and that this ignorance has been the root of our suffering. \emph{Gurus} and the scriptures play a vital role in guiding us out of the shadows of our own minds and into the light of clarity and wisdom.
\end{inspiration}
\newpage

\slcol{\Index{एवमेतद्यथात्थ} त्वमात्मानं परमेश्वर । \\
द्रष्टुमिच्छामि ते रूपमैश्वरं पुरुषोत्तम ॥ ११.३ ॥ }
\translate{ēvamētadyathāttha tvamātmānaṁ paramēśvara | \\
draṣṭumicchāmi tē rūpamaiśvaraṁ puruṣōttama || 11.3 || }
\cquote{O Supreme Lord, indeed just as You have declared Yourself to be, I long to behold Your divine form, O Puruṣottama (Supreme Person).}
\slcol{\Index{मन्यसे यदि तच्छक्यं} मया द्रष्टुमिति प्रभो | \\
योगेश्वर ततो मे त्वं दर्शयात्मानमव्ययम् ॥ ११.४ ॥ }
\translate{manyasē yadi tacchakyaṁ mayā draṣṭumiti prabhō | 
yōgēśvara tatō mē tvaṁ darśayātmānamavyayam || 11.4 || }
\cquote{O Lord, if You think it possible for me to behold it, then, O Lord of \emph{Yoga}(Krishna), reveal to me Your imperishable Self-form.}
\slcol{श्रीभगवानुवाच —\\
\Index{पश्य मे पार्थ} रूपाणि शतशोऽथ सहस्रशः | \\
नानाविधानि दिव्यानि नानावर्णाकृतीनि च ॥ ११.५ ॥ }
\translate{śrībhagavānuvāca —\\
paśya mē pārtha rūpāṇi śataśō:'tha sahasraśaḥ | \\
nānāvidhāni divyāni nānāvarṇākr̥tīni ca || 11.5 || }
\cquote{Sri Bhagavān said: Behold, O Pārtha (son of Pṛthā, Arjuna), My forms by hundreds and thousands of diverse kinds and divine, of various colours and shapes.}
\slcol{\Index{पश्यादित्यान्वसून्रुद्रानश्विनौ} मरुतस्तथा । \\
बहून्यदृष्टपूर्वाणि पश्याश्चर्याणि भारत ॥ ११.६ ॥ }
\translate{paśyādityānvasūnrudrānaśvinau marutastathā | \\
bahūnyadr̥ṣṭapūrvāṇi paśyāścaryāṇi bhārata || 11.6 || }
\cquote{Behold the \emph{Ādityas}, the \emph{Vasus}, the \emph{Rudras}, the twin \emph{Aśvins}, and the \emph{Maruts} as well. Behold, O Bhārata (Arjuna), many wondrous forms never seen before.}

\newpage
\begin{mananam}{\mananamfont {Mananam Śloka - 3, 4}}
\normalfont\mananamtext Am I able to trust those things that I cannot perceive or experience directly? When I was a child, what led me to trust the words of my parents and elders? As I grew into adulthood, what made me place faith in the knowledge provided by the education system?
What is the value of personally experiencing things which I previously only understood in theory? How can such “personal knowing” bring transformation? 
Is there a distinction between perception and experience? From a spiritual perspective, can one experience things beyond the senses? If so, how is it possible to experience without relying on the faculties of perception?
\end{mananam}
\sreflect
\begin{inspiration}{\inspirationTitleFont{Inspiration}}
\normalfont\mananamtext We are inclined to believe only in what we can perceive through our senses. When our understanding remains bound to sensory input, the phrase "seeing is believing" becomes the basis of our learning. \emph{Bhakti Yoga}, the path of devotion, promises a divine vision—an experience of God, the Supreme Being—perceived as vividly as the tangible world around us. This profound experience convinces the sense-dependent mind to accept a reality that transcends ordinary perception. Such a spiritual revelation is attainable only by a \emph{sādhaka} whose inner faculties have been finely tuned and purified.
\end{inspiration}
\newpage

\slcol{\Index{इहैकस्थं जगत्कृत्स्नं} पश्याद्य सचराचरम् । \\
मम देहे गुडाकेश यच्चान्यद्द्रष्टुमिच्छसि ॥ ११.७ ॥ }
\translate{ihaikasthaṁ jagatkr̥tsnaṁ paśyādya sacarācaram | \\
mama dēhē guḍākēśa yaccānyaddraṣṭumicchasi || 11.7 || }
\cquote{Behold now, O Guḍākeśa (Conqueror of sleep, Arjuna), in My body gathered together in One the entire \emph{jagat} (universe), both moving and unmoving, and also whatever else you desire to see.}
\slcol{\Index{न तु मां शक्यसे} द्रष्टुमनेनैव स्वचक्षुषा । \\
दिव्यं ददामि ते चक्षुः पश्य मे योगमैश्वरम् ॥ ११.८ ॥ }
\translate{na tu māṁ śakyasē draṣṭumanēnaiva svacakṣuṣā | \\
divyaṁ dadāmi tē cakṣuḥ paśya mē yōgamaiśvaram || 11.8 || }
\cquote{However, you cannot behold Me with your own eyes. I grant you the \emph{divya cakṣuḥ} (divine sight); behold My divine power of \emph{yoga}.}
\slcol{सञ्जय उवाच —\\
\Index{एवमुक्त्वा ततो} राजन्महायोगेश्वरो हरिः । \\
दर्शयामास पार्थाय परमं रूपमैश्वरम् ॥ ११.९ ॥ }
\translate{sañjaya uvāca —\\
ēvamuktvā tatō rājanmahāyōgēśvarō hariḥ | \\
darśayāmāsa pārthāya paramaṁ rūpamaiśvaram || 11.9 || }
\translate{Sañjaya said: Having spoken thus, O King (Dhṛtarāṣṭra), Hari (Krishna), the great Lord of \emph{Yoga}, revealed to Pārtha (Arjuna) His supreme, lordly form.}
\slcol{\Index{अनेकवक्त्रनयन}मनेकाद्भुतदर्शनम् । \\
अनेकदिव्याभरणं दिव्यानेकोद्यतायुधम् ॥ ११.१० ॥ }
\translate{anēkavaktranayanamanēkādbhutadarśanam | \\
anēkadivyābharaṇaṁ divyānēkōdyatāyudham || 11.10 || }
\cquote{(Lord’s form was revealed as) Having countless mouths and eyes, displaying countless wondrous visions, adorned with countless divine ornaments, and bearing countless divine uplifted weapons.}

\newpage
\begin{mananam}{\mananamfont {Mananam Śloka - 5, 7}}
\normalfont\mananamtext What does it signify for the One to manifest in many forms? Does this Being shift between forms, akin to an actor changing costumes? Do the various colours appear together, side by side, like the hues of a rainbow? Or is it a simultaneous expression of the same Being, manifesting in different forms and colours at once? Are these forms and colours confined to a specific space or time, or do they transcend such limitations, existing beyond the boundaries of space and time?
\end{mananam}
\sreflect
\begin{inspiration}{\inspirationTitleFont{Inspiration}}
\normalfont\mananamtext Forms, shapes and colours are the attributes through which we can distinguish one object from another. The Divine Being is free from all of these distinguishing attributes, and yet can manifest through any of them. Such an experience of seeing infinite forms manifested in One Being is the \emph{viśvarūpa darshan}, the vision of the universal form! Such a vision is not a physical phenomenon experienced through the senses.
\end{inspiration}
\newpage

\slcol{\Index{दिव्यमाल्याम्बरधरं} दिव्यगन्धानुलेपनम् । \\
सर्वाश्चर्यमयं देवमनन्तं विश्वतोमुखम् ॥ ११.११ ॥ }
\translate{divyamālyāmbaradharaṁ divyagandhānulēpanam |\\ 
sarvāścaryamayaṁ dēvamanantaṁ viśvatōmukham || 11.11 || }
\cquote{Adorned with divine garlands and garments, anointed with divine fragrances, full of all wonders, infinite, and facing everywhere, the shining Lord appeared.}
\slcol{\Index{दिवि सूर्यसहस्रस्य} भवेद्युगपदुत्थिता । \\
यदि भाः सदृशी सा स्याद्भासस्तस्य महात्मनः ॥ ११.१२ ॥ }
\translate{divi sūryasahasrasya bhavēdyugapadutthitā | \\
yadi bhāḥ sadr̥śī sā syādbhāsastasya mahātmanaḥ || 11.12 ||}
\cquote{If a thousand suns were to rise simultaneously in the sky, their splendour might resemble the radiance of that great Self (Supreme One).}
\slcol{\Index{तत्रैकस्थं जगत्कृत्स्नं} प्रविभक्तमनेकधा । \\
अपश्यद्देवदेवस्य शरीरे पाण्डवस्तदा ॥ ११.१३ ॥ }
\translate{tatraikasthaṁ jagatkr̥tsnaṁ pravibhaktamanēkadhā | \\
apaśyaddēvadēvasya śarīrē pāṇḍavastadā || 11.13 || }
\cquote{Then Pāṇḍava (Arjuna) beheld in the body of the Lord of lords (Krishna) the entire \emph{jagat} (universe), gathered together, yet divided into many, in that One.}
\slcol{\Index{ततः स विस्मयाविष्टो} हृष्टरोमा धनञ्जयः । \\
प्रणम्य शिरसा देवं कृताञ्जलिरभाषत ॥ ११.१४ ॥ }
\translate{tataḥ sa vismayāviṣṭō hr̥ṣṭarōmā dhanañjayaḥ | \\
praṇamya śirasā dēvaṁ kr̥tāñjalirabhāṣata || 11.14 || }\footnote{In the narrative of the Bhagavad Gītā, this moment of divine revelation is mirrored in the shift of focus to the narrator, Sanjaya. As an observer from afar, Sanjaya embodies the impartial witnessing awareness within each individual, representing the serene illumination that comes with inner realization.}
\cquote{Then Dhanañjaya, filled with wonder and with his hair standing on end, bowed his head to the shining Lord and spoke with joined palms.}


\newpage
\begin{mananam}{\mananamfont {Mananam Śloka - 10, 11, 12, 13, 14}}
\normalfont\mananamtext What is the most awe-inspiring experience of grandeur I have ever had—whether it was witnessing towering mountains, the vastness of the ocean, the stars scattered across the night sky, the boundless open fields, or something else of great splendour? How did such an experience influence my mind and emotions? 
Was I able, in that moment, to articulate the essence of what I was experiencing? Did I realize that my stream of thoughts momentarily ceased, allowing for a stillness within? When my mind resumed its activity to comprehend the nature of that experience, did I discover a renewed perspective that allowed me to embrace the experience more fully? Can I consciously cultivate this awareness and observe this process in my future encounters with life’s profound moments?
\end{mananam}
\sreflect
\begin{inspiration}{\inspirationTitleFont{Inspiration}}
\normalfont\mananamtext When beholding the most exquisite or magnificent aspects of nature, there is a fleeting pause in the stream of the mind. The mind, typically caught in its continuous cycle of judgment and evaluation—except during sleep—steps back in awe at the resplendence of the outer experience and adopts the role of a spectator. In this brief pause, the essence of one’s consciousness emerges, unclouded by the filters of the mind. Thus, such extraordinary human experiences indirectly guide one back to the ever-present, radiant consciousness within. It is at this point that the mind humbly surrenders and bows to the Divine essence within, from which all its strength and capabilities arise.
\end{inspiration}
\newpage

\slcol{अर्जुन उवाच —\\
\Index{पश्यामि देवांस्तव} देव देहे \\\hspace*{0.5cm}सर्वांस्तथा भूतविशेषसङ्घान् । \\
ब्रह्माणमीशं कमलासनस्थ-\\\hspace*{0.5cm}मृषींश्च सर्वानुरगांश्च दिव्यान् ॥ ११.१५ ॥ }
\translate{arjuna uvāca —\\
paśyāmi dēvāṁstava dēva dēhē \\
\hspace*{0.5cm}sarvāṁstathā bhūtaviśēṣasaṅghān | \\
brahmāṇamīśaṁ kamalāsanasthamr̥ṣīṁśca \\
\hspace*{0.5cm}sarvānuragāṁśca divyān || 11.15 || }
\cquote{Arjuna said: O Lord, I behold within Your body all the \emph{devas} (gods), and the hosts of beings; Brahmā, the Lord of creation, seated on the lotus, all the sages, and the divine serpents.}
\slcol{\Index{अनेकबाहूदरवक्त्रनेत्रं} \\\hspace*{0.5cm}पश्यामि त्वां सर्वतोऽनन्तरूपम् |\\
नान्तं न मध्यं न पुनस्तवादिं \\\hspace*{0.5cm}पश्यामि विश्वेश्वर विश्वरूप ॥ ११.१६ ॥ }
\translate{anekabāhūdaravaktranetraṁ paśyāmi \\
\hspace*{0.5cm}paśyāmi tvāṁ sarvato'nantarūpam |\\
nāntaṁ na madhyaṁ na punastavādiṁ \\
\hspace*{0.5cm}paśyāmi viśveśvara viśvarūpa || 11.16 || }
\cquote{I behold You everywhere with countless arms, bellies, mouths, and eyes, of infinite forms. I see neither end, nor middle, nor beginning of You, O Viśveśvara (Lord of Universe, Krishna), in Your \emph{viśvarūpa} (universal form).}

\newpage
\begin{mananam}{\mananamfont {Mananam Śloka - 15}}
\normalfont\mananamtext Can I accept the concept that, just as an organization operates under the leadership of various managers and executives, this creation is similarly governed by multiple gods and divine beings? If the forces of nature are directed and managed by these divine entities, what ensures their harmonious collaboration? Unlike an organization, where the leader is subject to being replaced or retiring, what distinguishes the Supreme Leader or controller of the cosmos? Is this ultimate God merely someone who possesses greater power than all others, or is there a deeper distinction?
\end{mananam}
\sreflect
\begin{inspiration}{\inspirationTitleFont{Inspiration}}
\normalfont\mananamtext Hindu cosmology is a richly detailed and intricate framework which harmoniously encompasses diverse perspectives, including monotheism, polytheism, pantheism, and panentheism. It offers representations of virtually every concept of God found in religions worldwide. Similarly, various creation narratives and the principles governing the forces of nature—both constructive and destructive—are thoughtfully considered and depicted within this cosmology. This openness provides individuals the freedom to embrace any concept of God or gods that aligns with their personal inclinations. For those willing to delve deeper, the Bhagavad Gītā and other Indian scriptures offer a comprehensive and unified cosmological vision that addresses all dimensions of creation and beyond.
\end{inspiration}
\newpage
\begin{mananam}{\mananamfont {Mananam Śloka - 16}}
\normalfont\mananamtext What is my understanding of the Supreme Being? Is it possible for anyone or anything to exist prior to That? Is it worthwhile to seek realization and knowledge of this Supreme Being? How might such understanding serve a purpose in my life? Could the Lord, revered as the most ancient and supreme of all, possess a purpose in upholding righteousness? What significance does the preservation of justice and order hold in maintaining the harmony and balance of the universe?
\end{mananam}
\sreflect
\begin{inspiration}{\inspirationTitleFont{Inspiration}}
\normalfont\mananamtext Only an imperishable Being can truly be the most ancient. Unlike humans, who often engage in struggles for supremacy and power, this eternal source of all existence remains unthreatened and secure, never fearing replacement. Every entity that arises from this Supreme Being carries its essence, representing a fragment of its infinite existence. The Lord, as the Creator and source of all beings, sustains and ensures the survival of all creation.
\end{inspiration}
\newpage

\slcol{\Index{किरीटिनं गदिनं} चक्रिणं च \\
\hspace*{0.5cm}तेजोराशिं सर्वतोदीप्तिमन्तम् । \\
पश्यामि त्वां दुर्निरीक्ष्यं समन्ताद्  \\
\hspace*{0.5cm}दीप्तानलार्कद्युतिमप्रमेयम् ॥ ११.१७ ॥ }
\translate{kirīṭinaṁ gadinaṁ cakriṇaṁ ca \\
\hspace*{0.5cm}tējōrāśiṁ sarvatōdīptimantam | \\
paśyāmi tvāṁ durnirīkṣyaṁ samantād\\
\hspace*{0.5cm}dīptānalārkadyutimapramēyam || 11.17 || }
\cquote{I behold You crowned, bearing a mace, holding a discus in hand, a mass of radiance blazing on all sides. I perceive, though difficult to behold, shining like burning fire and the sun, beyond all measure.}
\slcol{\Index{त्वमक्षरं परमं} वेदितव्यं \\
\hspace*{0.5cm}त्वमस्य विश्वस्य परं निधानम् । \\
त्वमव्ययः शाश्वतधर्मगोप्ता \\
\hspace*{0.5cm}सनातनस्त्वं पुरुषो मतो मे ॥ ११.१८ ॥ }
\translate{tvamakṣaraṁ paramaṁ vēditavyaṁ \\
\hspace*{0.5cm}tvamasya viśvasya paraṁ nidhānam | \\
tvamavyayaḥ śāśvatadharmagōptā \\
\hspace*{0.5cm}sanātanastvaṁ puruṣō matō mē || 11.18 || }
\cquote{You are the imperishable, the supreme reality to be known. You are the ultimate refuge of this universe. You are undecaying, the protector of the eternal \emph{dharma}, the eternal \emph{Puruṣa}—thus I regard You.}
\slcol{\Index{अनादिमध्यान्तम}नन्तवीर्यम्\\
\hspace*{0.5cm}अनन्तबाहुं शशिसूर्यनेत्रम् । \\
पश्यामि त्वां दीप्तहुताशवक्त्रं \\
\hspace*{0.5cm}स्वतेजसा विश्वमिदं तपन्तम् ॥ ११.१९ ॥ }
\translate{Anādimadhyāntamanantavīryam\\
\hspace*{0.5cm}anantabāhuṁ śaśisūryanētram | \\
paśyāmi tvāṁ dīptahutāśavaktraṁ\\
\hspace*{0.5cm}svatējasā viśvamidaṁ tapantam || 11.19 || }
\cquote{I behold You without beginning, middle, or end, of infinite power, with infinite arms, with the moon and sun as eyes, with blazing fire-like mouths, scorching this universe with Your own radiance.}
\slcol{\Index{द्यावापृथिव्योरिदमन्तरं} हि \\
\hspace*{0.5cm}व्याप्तं त्वयैकेन दिशश्च सर्वाः । \\
दृष्ट्वाद्भुतं रूपमिदं तवोग्रं \\
\hspace*{0.5cm}लोकत्रयं प्रव्यथितं महात्मन् ॥ ११.२० ॥ }
\translate{dyāvāpr̥thivyōridamantaraṁ hi \\
\hspace*{0.5cm}vyāptaṁ tvayaikēna diśaśca sarvāḥ | \\
dr̥ṣṭvādbhutaṁ rūpamidaṁ tavōgraṁ \\
\hspace*{0.5cm}lōkatrayaṁ pravyathitaṁ mahātman || 11.20 || }
\cquote{The space between heaven and earth, and all the directions, are pervaded by You alone. Beholding this wondrous and terrible form of Yours, O Great One, the three worlds are greatly agitated.}
\slcol{\Index{अमी हि त्वा} सुरसङ्घा विशन्ति \\
\hspace*{0.5cm}केचिद्भीताः प्राञ्जलयो गृणन्ति । \\
स्वस्तीत्युक्त्वा महर्षिसिद्धसङ्घाः \\
\hspace*{0.5cm}स्तुवन्ति त्वां स्तुतिभिः पुष्कलाभिः ॥११.२१॥ }
\translate{amī hi tvāṁ surasaṅghā viśanti \\
\hspace*{0.5cm}kēcidbhītāḥ prāñjalayō gr̥ṇanti | \\
svastītyuktvā maharṣisiddhasaṅghāḥ \\
\hspace*{0.5cm}stuvanti tvāṁ stutibhiḥ puṣkalābhiḥ || 11.21 || }
\cquote{Indeed, these hosts of gods enter into You; some, fearful, praise You with joined palms. The great sages and perfected beings, saying ‘May there be peace,’ glorify You with abundant hymns of praise.}
\slcol{\Index{रुद्रादित्या वसवो} ये च साध्या \\
\hspace*{0.5cm}विश्वेऽश्विनौ मरुतश्चोष्मपाश्च । \\\
गन्धर्वयक्षासुरसिद्धसङ्घा \\
\hspace*{0.5cm}वीक्षन्ते त्वां विस्मिताश्चैव सर्वे ॥ ११.२२ ॥ }
\translate{rudrādityā vasavō yē ca sādhyā \\
\hspace*{0.5cm}viśvē:'śvinau marutaścōṣmapāśca | \\
gandharvayakṣāsurasiddhasaṅghā \\
\hspace*{0.5cm}vīkṣantē tvāṁ vismitāścaiva sarvē || 11.22 || }
\cquote{The \emph{Rudras}, \emph{Ādityas}, \emph{Vasus}, and \emph{Sādhyas}; the \emph{Viśvedevas}, the \emph{Aśvins}, the \emph{Maruts}, and the guardians of fire; the hosts of \emph{Gandharvas}, \emph{Yakṣas}, \emph{Asuras}, and \emph{Siddhas}—all gaze upon You, astonished indeed.}
\slcol{\Index{रूपं महत्ते बहुवक्त्रनेत्रं} \\
\hspace*{0.5cm}महाबाहो बहुबाहूरुपादम् । \\
बहूदरं बहुदंष्ट्राकरालं \\
\hspace*{0.5cm}दृष्ट्वा लोकाः प्रव्यथितास्तथाहम् ॥ ११.२३ ॥ }
\translate{rūpaṁ mahattē bahuvaktranētraṁ \\
\hspace*{0.5cm}mahābāhō bahubāhūrupādam | \\
bahūdaraṁ bahudaṁṣṭrākarālaṁ \\
\hspace*{0.5cm}dr̥ṣṭvā lōkāḥ pravyathitāstathāham || 11.23 || }
\cquote{O mighty-armed, I behold Your great form with many mouths and eyes, with many arms, thighs, and feet, with many bellies, and terrifying with many fangs. Seeing this, the worlds are agitated, and so am I.}
\slcol{\Index{नभःस्पृशं दीप्तमनेकवर्णं} \\
\hspace*{0.5cm}व्यात्ताननं दीप्तविशालनेत्रम् । \\
दृष्ट्वा हि त्वां प्रव्यथितान्तरात्मा \\
\hspace*{0.5cm}धृतिं न विन्दामि शमं च विष्णो ॥ ११.२४ ॥ }
\translate{nabhaḥspr̥śaṁ dīptamanēkavarṇaṁ \\
\hspace*{0.5cm}vyāttānanaṁ dīptaviśālanētram | \\
dr̥ṣṭvā hi tvāṁ pravyathitāntarātmā \\
\hspace*{0.5cm}dhr̥tiṁ na vindāmi śamaṁ ca viṣṇō || 11.24 || }
\cquote{Touching the sky, blazing in many colours, with gaping mouths and vast blazing eyes—beholding You, my inner self is agitated. I find no composure, no peace, O Viṣṇu.}
\slcol{\Index{दंष्ट्राकरालानि च} ते मुखानि \\
\hspace*{0.5cm}दृष्ट्वैव कालानलसंनिभानि । \\
दिशो न जाने न लभे च शर्म \\
\hspace*{0.5cm}प्रसीद देवेश जगन्निवास ॥ ११.२५ ॥ }
\translate{daṁṣṭrākarālāni ca tē mukhāni \\
\hspace*{0.5cm}dr̥ṣṭvaiva kālānalasaṁnibhāni | \\
diśō na jānē na labhē ca śarma \\
\hspace*{0.5cm}prasīda dēvēśa jagannivāsa || 11.25 || }
\cquote{Beholding Your mouths, terrible with fangs, resembling the fire of Time (of destruction), I have lost sense of direction, nor do I find peace. Be gracious, O Lord of lords, Abode of the universe.”}
\slcol{\Index{अमी च त्वां} धृतराष्ट्रस्य पुत्राः \\
\hspace*{0.5cm}सर्वे सहैवावनिपालसङ्घैः । \\
भीष्मो द्रोणः सूतपुत्रस्तथासौ \\
\hspace*{0.5cm}सहास्मदीयैरपि योधमुख्यैः ॥ ११.२६ ॥ }
\translate{amī ca tvāṁ dhr̥tarāṣṭrasya putrāḥ \\
\hspace*{0.5cm}sarvē sahaivāvanipālasaṅghaiḥ | \\
bhīṣmō drōṇaḥ sūtaputrastathāsau \\
\hspace*{0.5cm}sahāsmadīyairapi yōdhamukhyaiḥ || 11.26 || }
\cquote{The sons of Dhṛtarāṣṭra, all together with the hosts of kings; Bhīṣma, Droṇa, and the son of charioteer (Karṇa) likewise, along with the foremost of our own warriors…}
\slcol{\Index{वक्त्राणि ते} त्वरमाणा विशन्ति \\
\hspace*{0.5cm}दंष्ट्राकरालानि भयानकानि । \\
केचिद्विलग्ना दशनान्तरेषु \\
\hspace*{0.5cm}सन्दृश्यन्ते चूर्णितैरुत्तमाङ्गैः ॥ ११.२७ ॥ }
\translate{vaktrāṇi tē tvaramāṇā viśanti \\
\hspace*{0.5cm}daṁṣṭrākarālāni bhayānakāni | \\
kēcidvilagnā daśanāntarēṣu \\
\hspace*{0.5cm}sandr̥śyantē cūrṇitairuttamāṅgaiḥ || 11.27 || }
\cquote{All these warriors rush into Your fierce mouths, terrible with fangs, frightening. Some are caught in between the gaps of the teeth, and their heads are seen crushed to pieces.}

\newpage
\begin{mananam}{\mananamfont {Mananam Śloka - 23, 24, 25}}
\normalfont\mananamtext Why does the sight of something immense or the presence of someone powerful instil fear in the heart? How is it that the vision of the Supreme Being—who is both worshipped and adored—can also evoke a sense of fear? Can this be compared to our relationships with certain authority figures in our lives, where we feel both love and a measure of fear toward them? What distinguishes the power wielded by the Lord, the most powerful being, from the authority exercised by human dictators over their people? 
\end{mananam}
\sreflect
\begin{inspiration}{\inspirationTitleFont{Inspiration}}
\normalfont\mananamtext Indian deities are often depicted not only as loving and endearing figures but also as powerful and, at times, fierce manifestations of the Supreme. The same Supreme Being who showers blessings upon devotees with boundless compassion is also mightier than any demon, capable of vanquishing them completely.
When we remain unaware of our relationship with the Infinite, we continue to identify ourselves as limited, finite beings. Fear arises from the perception of duality and separation—of there being an “other.” However, the realization of our true essence and oneness with the Supreme allows love, rather than fear, to become the guiding and dominant force that governs the heart.
\end{inspiration}
\newpage

\slcol{\Index{यथा नदीनां} बहवोऽम्बुवेगाः \\
\hspace*{0.5cm}समुद्रमेवाभिमुखा द्रवन्ति । \\
तथा तवामी नरलोकवीरा \\
\hspace*{0.5cm}विशन्ति वक्त्राण्यभिविज्वलन्ति ॥ ११.२८ ॥ }
\translate{yathā nadīnāṁ bahavō:'mbuvēgāḥ \\
\hspace*{0.5cm}samudramēvābhimukhā dravanti | \\
tathā tavāmī naralōkavīrā \\
\hspace*{0.5cm}viśanti vaktrāṇyabhivijvalanti || 11.28 || }
\cquote{Just as many torrents of rivers rush towards the ocean, so too these heroes of the human world enter into Your blazing mouths.}
\slcol{\Index{यथा प्रदीप्तं} ज्वलनं पतङ्गाः\\
\hspace*{0.5cm}विशन्ति नाशाय समृद्धवेगाः । \\
तथैव नाशाय विशन्ति लोकाः \\
\hspace*{0.5cm}तवापि वक्त्राणि समृद्धवेगाः ॥ ११.२९ ॥ }
\translate{yathā pradīptaṁ jvalanaṁ pataṅgāḥ\\
\hspace*{0.5cm}viśanti nāśāya samr̥ddhavēgāḥ | \\
tathaiva nāśāya viśanti lōkāḥ \\
\hspace*{0.5cm}tavāpi vaktrāṇi samr̥ddhavēgāḥ || 11.29 || }
\cquote{Just as moths rush swiftly into a blazing fire to their destruction, so too these beings rush swiftly into Your mouths for their destruction.}
\slcol{\Index{लेलिह्यसे ग्रसमानः} समन्ताद्\\
\hspace*{0.5cm}लोकान्समग्रान्वदनैर्ज्वलद्भिः । \\
तेजोभिरापूर्य जगत्समग्रं \\
\hspace*{0.5cm}भासस्तवोग्राः प्रतपन्ति विष्णो ॥ ११.३० ॥ }
\translate{lēlihyasē grasamānaḥ samantād\\
\hspace*{0.5cm}lōkānsamagrānvadanairjvaladbhiḥ | \\
tējōbhirāpūrya jagatsamagraṁ \\
\hspace*{0.5cm}bhāsastavōgrāḥ pratapanti viṣṇō || 11.30 || }
\cquote{Devouring all the worlds from every side with Your blazing mouths, You lick them in delight. Your fierce brilliance fills the entire universe with radiance and scorches everything, O Viṣṇu.}
\slcol{\Index{आख्याहि मे को} भवानुग्ररूप: \\
\hspace*{0.5cm}नमोऽस्तु ते देववर प्रसीद । \\
विज्ञातुमिच्छामि भवन्तमाद्यं \\
\hspace*{0.5cm}न हि प्रजानामि तव प्रवृत्तिम् ॥ ११.३१ ॥ }
\translate{ākhyāhi mē kō bhavānugrarūpō \\
\hspace*{0.5cm}namō:'stu tē dēvavara prasīda | \\
vijñātumicchāmi bhavantamādyaṁ \\
\hspace*{0.5cm}na hi prajānāmi tava pravr̥ttim || 11.31 || }
\cquote{Tell me, who are You in this terrible form? Salutations to You, O Supreme Deva, be gracious. I wish to know You, the primal Being, for I do not understand Your purpose.}
\slcol{श्रीभगवानुवाच —\\
\Index{कालोऽस्मि} लोकक्षयकृत्प्रवृद्धो\\
\hspace*{0.5cm}लोकान्समाहर्तुमिह प्रवृत्तः । \\
ऋतेऽपि त्वा न भविष्यन्ति सर्वे \\
\hspace*{0.5cm}येऽवस्थिताः प्रत्यनीकेषु योधाः ॥ ११.३२ ॥ }
\translate{śrībhagavānuvāca —\\
kālō:'smi lōkakṣayakr̥tpravr̥ddhō\\
\hspace*{0.5cm}lōkānsamāhartumiha pravr̥ttaḥ | \\ 
r̥tē:'pi tvā na bhaviṣyanti sarvē \\
\hspace*{0.5cm}yē:'vasthitāḥ pratyanīkēṣu yōdhāḥ || 11.32 || }
\cquote{Sri Bhagavān said: I am Time, the great destroyer of worlds, engaged in annihilating all the worlds. Even without you, all the warriors arrayed in the opposing ranks shall cease to exist.}

\newpage
\begin{mananam}{\mananamfont {Mananam Śloka - 28, 29, 30, 31}}
\normalfont\mananamtext How can destruction be considered one of the activities of the Supreme? How is it that the same Being who creates is also responsible for destruction? How can these seemingly opposing functions coexist within a single entity?
What positive aspects can be found in the act of destruction? Could the destructive force present in nature also serve to eliminate bad habits and negativity within my own inner self?
\end{mananam}
\sreflect
\begin{inspiration}{\inspirationTitleFont{Inspiration}}
\normalfont\mananamtext Within the Supreme Being lies the origin and dissolution of all creation. The world, along with every individual, follows the inevitable cycle of change—from birth to growth to decay and eventually, to death. Embracing this constant change becomes a pathway to inner peace and personal growth.
Destruction, often perceived as negative, is a vital aspect of nature that preserves its dynamic freshness. Through this continual change, life offers all individuals the opportunity to begin anew and make the most of their lives. It is this very force of destruction that enables spiritual resurgence, purification, and ultimately, liberation.
\end{inspiration}
\newpage

\slcol{\Index{तस्मात्त्वमुत्तिष्ठ} यशो लभस्व \\
\hspace*{0.5cm}जित्वा शत्रून्भुङ्क्ष्व राज्यं समृद्धम् । \\
मयैवैते निहताः पूर्वमेव \\
\hspace*{0.5cm}निमित्तमात्रं भव सव्यसाचिन् ॥ ११.३३ ॥ }
\translate{tasmāttvamuttiṣṭha yaśō labhasva \\
\hspace*{0.5cm}jitvā śatrūnbhuṅkṣva rājyaṁ samr̥ddham | \\
mayaivaitē nihatāḥ pūrvamēva \\
\hspace*{0.5cm}nimittamātraṁ bhava savyasācin || 11.33 || }
\cquote{Therefore, arise and claim your glory. Defeat your enemies and enjoy a prosperous kingdom. These warriors have already been slain by Me alone. Become merely an instrument (of Mine), O left-handed archer (Arjuna).}
\slcol{\Index{द्रोणं च भीष्मं} च जयद्रथं च \\
\hspace*{0.5cm}कर्णं तथान्यानपि योधवीरान् । \\
मया हतांस्त्वं जहि मा व्यथिष्ठा \\
\hspace*{0.5cm}युध्यस्व जेतासि रणे सपत्नान् ॥ ११.३४ ॥ }
\translate{drōṇaṁ ca bhīṣmaṁ ca jayadrathaṁ ca \\
\hspace*{0.5cm}karṇaṁ tathānyānapi yōdhavīrān | \\
mayā hatāṁstvaṁ jahi mā vyathiṣṭhā \\
\hspace*{0.5cm}yudhyasva jētāsi raṇē sapatnān || 11.34 || }
\cquote{Droṇa, Bhīṣma, Jayadratha, Karṇa, and many other heroic warriors have already been slain by Me. Therefore, kill them; do not be disturbed (or conflicted within). Fight, and you shall conquer your enemies in battle.}

\newpage
\begin{mananam}{\mananamfont {Mananam Śloka - 32, 33}}
\normalfont\mananamtext How can Time be understood as a force of destruction? In what manner does the Lord exercise control over Time? To what extent are our lives already predestined? 
How does the Lord bring about accomplishments independently of our cooperation? Are there not powerful individuals who seem indispensable, without whom the course of events would change? How can adopting the perspective that the Lord is the ultimate doer while we are merely instruments be beneficial to a devotee?
\end{mananam}
\sreflect
\begin{inspiration}{\inspirationTitleFont{Inspiration}}
\normalfont\mananamtext Change is the only constant force in nature. It governs the cycle of seasons and ensures that every living being born will inevitably face death. Neither birth nor death is within our control. This profound truth serves as a humbling reminder, even to the most accomplished and powerful individuals. When we truly internalize this reality and surrender ourselves completely to the Divine, we become worthy instruments of His will. Genuine success and lasting fame come to those who recognize that God is the ultimate doer, and we are but vessels for His work.
\end{inspiration}
\newpage

\slcol{सञ्जय उवाच —\\
\Index{एतच्छ्रुत्वा वचनं} केशवस्य \\\hspace*{0.5cm}कृताञ्जलिर्वेपमानः किरीटी । \\
नमस्कृत्वा भूय एवाह कृष्णं \\\hspace*{0.5cm}सगद्गदं भीतभीतः प्रणम्य ॥ ११.३५ ॥ }
\translate{sañjaya uvāca —\\
ētacchrutvā vacanaṁ kēśavasya \\
\hspace*{0.5cm}kr̥tāñjalirvēpamānaḥ kirīṭī | \\
namaskr̥tvā bhūya ēvāha kr̥ṣṇaṁ \\
\hspace*{0.5cm}sagadgadaṁ bhītabhītaḥ praṇamya || 11.35 || }
\cquote{Sañjaya said: Having heard these words of Keśava (Krishna), the crowned warrior (Arjuna), trembling and with folded palms, bowed down again. With a faltering voice, overcome with fear, he spoke.}
\slcol{अर्जुन उवाच —\\
\Index{स्थाने हृषीकेश} तव प्रकीर्त्या \\
\hspace*{0.5cm}जगत्प्रहृष्यत्यनुरज्यते च । \\
रक्षांसि भीतानि दिशो द्रवन्ति \\
\hspace*{0.5cm}सर्वे नमस्यन्ति च सिद्धसङ्घाः ॥ ११.३६ ॥ }
\translate{arjuna uvāca —\\
sthānē hr̥ṣīkēśa tava prakīrtyā \\
\hspace*{0.5cm}jagatprahr̥ṣyatyanurajyatē ca | \\
rakṣāṁsi bhītāni diśō dravanti \\
\hspace*{0.5cm}sarvē namasyanti ca siddhasaṅghāḥ || 11.36 || }
\cquote{Arjuna said: Rightly, O Hṛṣīkeśa (Krishna), by Your praise the world rejoices and becomes devoted. The demons, terrified, flee in all directions, while all the hosts of perfected beings (\emph{siddhas}) bow down before You.}
\slcol{\Index{कस्माच्च ते न} नमेरन्महात्मन्\\
\hspace*{0.5cm}गरीयसे ब्रह्मणोऽप्यादिकर्त्रे । \\
अनन्त देवेश जगन्निवास \\
\hspace*{0.5cm}त्वमक्षरं सदसत्तत्परं यत् ॥ ११.३७ ॥ }
\translate{kasmācca tē na namēranmahātman\\
\hspace*{0.5cm}garīyasē brahmaṇō:'pyādikartrē | \\
ananta dēvēśa jagannivāsa \\
\hspace*{0.5cm}tvamakṣaraṁ sadasattatparaṁ yat || 11.37 || }
\cquote{Why should they not bow to You, O great souled-One? You are greater even than Brahmā, the primal creator. O Infinite One, O Lord of the gods, O Abode of the universe, You are the imperishable, beyond both the manifest and the unmanifest.}
\slcol{\Index{त्वमादिदेवः पुरुषः} पुराणः\\
\hspace*{0.5cm}त्वमस्य विश्वस्य परं निधानम् । \\
वेत्तासि वेद्यं च परं च धाम \\
\hspace*{0.5cm}त्वया ततं विश्वमनन्तरूप ॥ ११.३८ ॥ }
\translate{tvamādidēvaḥ puruṣaḥ purāṇaḥ\\
\hspace*{0.5cm}tvamasya viśvasya paraṁ nidhānam | \\
vēttāsi vēdyaṁ ca paraṁ ca dhāma \\
\hspace*{0.5cm}tvayā tataṁ viśvamanantarūpa || 11.38 || }
\cquote{You are the primal \emph{Deva} (God), the ancient Person. You are the supreme resting place of this universe. You are the knower, the knowable, and the supreme abode. By You, the infinite-formed, the universe is pervaded.}
\slcol{\Index{वायुर्यमोऽग्निर्वरुणः} शशाङ्कः \\
\hspace*{0.5cm}प्रजापतिस्त्वं प्रपितामहश्च । \\
नमो नमस्तेऽस्तु सहस्रकृत्वः \\
\hspace*{0.5cm}पुनश्च भूयोऽपि नमो नमस्ते ॥ ११.३९ ॥ }
\translate{vāyuryamō:'gnirvaruṇaḥ śaśāṅkaḥ \\
\hspace*{0.5cm}prajāpatistvaṁ prapitāmahaśca | \\
namō namastē:'stu sahasrakr̥tvaḥ \\
\hspace*{0.5cm}punaśca bhūyō:'pi namō namastē || 11.39 || }
\cquote{You are Vāyu, Yama, Agni, Varuṇa, the moon; You are Prajāpati and the great-grandfather of all. Salutations to You a thousand times, and again and again, salutations to You.}

\newpage
\begin{mananam}{\mananamfont {Mananam Śloka - 37, 38, 39, 40}}
\normalfont\mananamtext Why are descriptions of the Divine often beyond the limits of human imagination? Why should we refrain from taking ultimate refuge in a human form or an institutional system? What is the significance of bowing to the Supreme Lord? How does bowing to the Lord differ from submitting to a human superior or an authority figure? Should the expression “bowing down in all directions” be understood literally or metaphorically?
\end{mananam}
\sreflect
\begin{inspiration}{\inspirationTitleFont{Inspiration}}
\normalfont\mananamtext The Supreme Being represents the subtlest of all conceptions, serving as a pointer toward that which transcends the boundaries of human thought. When we reach the limits of our mental constructs, we often feel frustration and exasperation. However, it is precisely in this state—beyond the grasp of the mind—that an opportunity arises. Scriptures use this seeming dead-end of the mind to guide the seeker toward that which lies beyond the mind and intellect. It is a state of profound humility, where one surrenders completely to the Source and the ultimate destination of one’s Being.
\end{inspiration}
\newpage

\slcol{\Index{नमः पुरस्तादथ} पृष्ठतस्ते \\
\hspace*{0.5cm}नमोऽस्तु ते सर्वत एव सर्व । \\
अनन्तवीर्यामितविक्रमस्त्वं \\
\hspace*{0.5cm}सर्वं समाप्नोषि ततोऽसि सर्वः ॥ ११.४० ॥ }
\translate{namaḥ purastādatha pr̥ṣṭhatastē \\
\hspace*{0.5cm}namō:'stu tē sarvata ēva sarva | \\
anantavīryāmitavikramastvaṁ \\
\hspace*{0.5cm}sarvaṁ samāpnōṣi tatō:'si sarvaḥ || 11.40 || }
\cquote{Salutations to You in front and behind, salutations to You from all sides, O All. You are of infinite power and immeasurable might. You pervade everything, therefore You are everything.}
\slcol{\Index{सखेति मत्वा} प्रसभं यदुक्तं \\\hspace*{0.5cm}हे कृष्ण हे यादव हे सखेति । \\
अजानता महिमानं तवेदं \\\hspace*{0.5cm}मया प्रमादात्प्रणयेन वापि ॥ ११.४१ ॥}
\translate{sakhēti matvā prasabhaṁ yaduktaṁ \\
\hspace*{0.5cm}hē kr̥ṣṇa hē yādava hē sakhēti | \\
ajānatā mahimānaṁ tavēdaṁ \\
\hspace*{0.5cm}mayā pramādātpraṇayēna vāpi || 11.41 ||}
\cquote{Thinking of You as a friend, whatever I have boldly said – O Krishna, O Yādava
 (descendant of Yadu clan), O friend–not knowing Your greatness, I have spoken either out of carelessness or out of affection.}
\slcol{\Index{यच्चावहासार्थमसत्कृतोऽसि} \\\hspace*{0.5cm}विहारशय्यासनभोजनेषु । \\
एकोऽथवाप्यच्युत तत्समक्षं \\\hspace*{0.5cm}तत्क्षामये त्वामहमप्रमेयम् ॥ ११.४२ ॥ }
\translate{yaccāvahāsārthamasatkr̥tō:'si \\
\hspace*{0.5cm}vihāraśayyāsanabhōjanēṣu | \\
ēkō:'thavāpyacyuta tatsamakṣaṁ \\
\hspace*{0.5cm}tatkṣāmayē tvāmahamapramēyam || 11.42 || }
\cquote{And for any disrespect I have shown You in jest – whether while playing, resting, sitting, or eating, whether alone, O Acyuta, or in the company of others – I humbly ask Your forgiveness, O immeasurable One (\emph{aprameyam}).}
\slcol{\Index{पितासि लोकस्य} चराचरस्य \\
\hspace*{0.5cm}त्वमस्य पूज्यश्च गुरुर्गरीयान् । \\
न त्वत्समोऽस्त्यभ्यधिकः कुतोऽन्यो \\
\hspace*{0.5cm}लोकत्रयेऽप्यप्रतिमप्रभाव ॥ ११.४३ ॥ }
\translate{pitāsi lōkasya carācarasya \\
\hspace*{0.5cm}tvamasya pūjyaśca gururgarīyān | \\
na tvatsamō:'styabhyadhikaḥ kutō:'nyō \\
\hspace*{0.5cm}lōkatrayē:'pyapratimaprabhāva || 11.43 || }
\cquote{You are the father of this world, of all that moves and all that is still. You are most worthy of worship and the greatest of \emph{gurus}. No one is equal to You; how then could there be anyone greater than You in all the worlds, even across the three realms, O One of incomparable power!}
\slcol{\Index{तस्मात्प्रणम्य} प्रणिधाय कायं \\\hspace*{0.5cm}प्रसादये त्वामहमीशमीड्यम् । \\
पितेव पुत्रस्य सखेव सख्युः \\\hspace*{0.5cm}प्रियः प्रियायार्हसि देव सोढुम् ॥११.४४॥}
\translate{tasmātpraṇamya praṇidhāya kāyaṁ \\
\hspace*{0.5cm}prasādayē tvāmahamīśamīḍyam | \\
pitēva putrasya sakhēva sakhyuḥ \\
\hspace*{0.5cm}priyaḥ priyāyārhasi dēva sōḍhum || 11.44 || }
\cquote{Therefore, bowing down and prostrating my body, I seek Your grace, O adorable Lord. As a father forgives his son, as a friend forgives his friend, as a beloved forgives the beloved—so should You forgive me, O Lord.}

\newpage
\begin{mananam}{\mananamfont {Mananam Śloka - 41, 42, 43, 44}}
\normalfont\mananamtext Have there been moments when you reflected on your past and realized that you had underestimated the significance of someone in your life? Can you connect with the sense of remorse for not expressing gratitude to those who truly deserved it? Have you ever engaged in trivialities with individuals who were worthy of your reverence and admiration?
Can you revisit your present life with a profound sense of appreciation and an open heart toward your \emph{Guru} and spiritual teachers? Are you offering them the respect they deserve and remaining receptive to their wisdom and guidance? If not, what shifts in your mindset and attitude could you adopt to cultivate consistent and intentional reverence toward them?
\end{mananam}
\sreflect
\begin{inspiration}{\inspirationTitleFont{Inspiration}}
\normalfont\mananamtext All our friends and well-wishers are, in truth, manifestations of the Divine, sent into our lives for our welfare and upliftment. Among these, the most profound manifestation is found in our spiritual teachers and guides, whose invaluable assistance helps us to navigate the challenging ocean of \emph{saṁsāra}. 
The presence of a \emph{guru} or spiritual teacher in our life is the greatest blessing. Yet, many people often overlook this gift, seeking a greater manifestation of the Divine in temples or amidst the grandeur of nature. Our faith and devotion tend to fluctuate, rising and falling like a roller coaster, driven by the ebb and flow of our emotions. However, when we learn to view life through the lens of wisdom rather than that of our passing feelings, then our hearts find stability and we embrace life with clarity.
\end{inspiration}
\newpage

\slcol{\Index{अदृष्टपूर्वं हृषितोऽस्मि} दृष्ट्वा \\\hspace*{0.5cm}भयेन च प्रव्यथितं मनो मे । \\
तदेव मे दर्शय देव रूपं \\\hspace*{0.5cm}प्रसीद देवेश जगन्निवास ॥ ११.४५ ॥ }
\translate{adr̥ṣṭapūrvaṁ hr̥ṣitō:'smi dr̥ṣṭvā \\
\hspace*{0.5cm}bhayēna ca pravyathitaṁ manō mē | \\
tadēva mē darśaya dēva rūpaṁ \\
\hspace*{0.5cm}prasīda dēvēśa jagannivāsa || 11.45 ||}
\cquote{I am delighted to see what has never been seen before, yet my mind is troubled with fear. Therefore, show me again Your gentle form. Be gracious, O Lord of gods, Abode of the universe.}
\slcol{\Index{किरीटिनं गदिनं} चक्रहस्तं \\\hspace*{0.5cm}इच्छामि त्वां द्रष्टुमहं तथैव । \\
तेनैव रूपेण चतुर्भुजेन \\\hspace*{0.5cm}सहस्रबाहो भव विश्वमूर्ते ॥ ११.४६ ॥ }
\translate{kirīṭinaṁ gadinaṁ cakrahastam \\
\hspace*{0.5cm}icchāmi tvāṁ draṣṭumahaṁ tathaiva | \\
tēnaiva rūpēṇa caturbhujēna \\
\hspace*{0.5cm}sahasrabāhō bhava viśvamūrtē || 11.46 || }
\cquote{I wish to see You as before, crowned, holding mace, and a discus in hand. O thousand-armed One of universal form, please appear in that very four-armed form of Yours.}
\slcol{श्रीभगवानुवाच —\\
\Index{मया प्रसन्नेन} तवार्जुनेदं \\\hspace*{0.5cm}रूपं परं दर्शितमात्मयोगात् । \\
तेजोमयं विश्वमनन्तमाद्यं \\\hspace*{0.5cm}यन्मे त्वदन्येन न दृष्टपूर्वम् ॥ ११.४७ ॥ }
\translate{śrībhagavānuvāca —\\
mayā prasannēna tavārjunēdaṁ \\
\hspace*{0.5cm}rūpaṁ paraṁ darśitamātmayōgāt | \\ 
tējōmayaṁ viśvamanantamādyaṁ \\
\hspace*{0.5cm}yanmē tvadanyēna na dr̥ṣṭapūrvam || 11.47 || }
\cquote{Sri Bhagavān said: By My grace, O Arjuna, this supreme form has been revealed to you through My own divine power. It is radiant, infinite, primal, encompassing the universe—a vision never before seen by anyone other than you. }
\slcol{\Index{न वेदयज्ञाध्ययनैर्न} दानैः \\\hspace*{0.5cm}न च क्रियाभिर्न तपोभिरुग्रैः । \\
एवंरूपः शक्य अहं नृलोके \\\hspace*{0.5cm}द्रष्टुं त्वदन्येन कुरुप्रवीर ॥ ११.४८ ॥ }
\translate{na vēdayajñādhyayanairna dānaiḥ \\
\hspace*{0.5cm}na ca kriyābhirna tapōbhirugraiḥ | \\
ēvaṁrūpaḥ śakya ahaṁ nr̥lōkē \\
\hspace*{0.5cm}draṣṭuṁ tvadanyēna kurupravīra || 11.48 || }
\cquote{Not by study of the \emph{vedas}, nor by sacrifices, nor by gifts, nor by rituals, nor by severe austerities can I be seen in this form in the world of men, by anyone other than you, O \emph{Kuru-pravīra} foremost of the \emph{Kurus}, Arjuna.}
\slcol{\Index{मा ते व्यथा मा च} विमूढभावः \\
\hspace*{0.5cm}दृष्ट्वा रूपं घोरमीदृङ्ममेदम् । \\
व्यपेतभीः प्रीतमनाः पुनस्त्वं \\
\hspace*{0.5cm}तदेव मे रूपमिदं प्रपश्य ॥ ११.४९ ॥ }
\translate{mā tē vyathā mā ca vimūḍhabhāvō \\
\hspace*{0.5cm}dr̥ṣṭvā rūpaṁ ghōramīdr̥ṅmamēdam | \\
vyapētabhīḥ prītamanāḥ punastvaṁ \\
\hspace*{0.5cm}tadēva mē rūpamidaṁ prapaśya || 11.49 || }
\cquote{Let there be no fear in you, nor bewilderment, on seeing this terrible form of Mine. Freed from fear, with a cheerful mind, behold again that very form of Mine.}
\slcol{सञ्जय उवाच —\\
\Index{इत्यर्जुनं} वासुदेवस्तथोक्त्वा \\
\hspace*{0.5cm}स्वकं रूपं दर्शयामास भूयः । \\
आश्वासयामास च भीतमेनं \\
\hspace*{0.5cm}भूत्वा पुनःसौम्यवपुर्महात्मा ॥ ११.५० ॥ }
\translate{sañjaya uvāca —\\
ityarjunaṁ vāsudēvastathōktvā \\
\hspace*{0.5cm}svakaṁ rūpaṁ darśayāmāsa bhūyaḥ | \\
āśvāsayāmāsa ca bhītamēnaṁ \\
\hspace*{0.5cm}bhūtvā punaḥsaumyavapurmahātmā || 11.50 || }
\cquote{Sañjaya said: Having spoken thus to Arjuna, Vāsudeva (Krishna) revealed His own form once again. The great soul, assuming His gentle and gracious appearance, comforted the terrified Arjuna.}
\slcol{अर्जुन उवाच —\\
\Index{दृष्ट्वेदं मानुषं रूपं} तव सौम्यं जनार्दन । \\
इदानीमस्मि संवृत्तः सचेताः प्रकृतिं गतः ॥ ११.५१ ॥ }
\translate{arjuna uvāca —\\
dr̥ṣṭvēdaṁ mānuṣaṁ rūpaṁ tava saumyaṁ janārdana | \\
idānīmasmi saṁvr̥ttaḥ sacētāḥ prakr̥tiṁ gataḥ || 11.51 || }
\cquote{Arjuna said: O Janārdana (Krishna), beholding Your gentle human form, I am now at peace. My mind has regained clarity, and I have returned to my natural state.}
\slcol{श्रीभगवानुवाच —\\
\Index{सुदुर्दर्शमिदं रूपं} दृष्टवानसि यन्मम । \\
देवा अप्यस्य रूपस्य नित्यं दर्शनकाङ्क्षिणः ॥ ११.५२ ॥ }
\translate{śrībhagavānuvāca —\\
sudurdarśamidaṁ rūpaṁ dr̥ṣṭavānasi yanmama | \\
dēvā apyasya rūpasya nityaṁ darśanakāṅkṣiṇaḥ || 11.52 || }
\cquote{Sri Bhagavān said: This form of Mine, which you have seen is exceedingly difficult to behold. Even the gods are ever longing to see this form.}

\newpage
\begin{mananam}{\mananamfont {Mananam Śloka - 50, 51}}
\normalfont\mananamtext How can a multi-form vision of the Divine, as perceived by Arjuna, be overwhelming? Why does connecting with the Lord in a particular human form bring comfort and peace to our hearts? How can the Lord’s presence offer solace, particularly when we are anxious or afraid?
Why does a devotee yearn for a personal relationship with God, especially in a human form? Is it possible to nurture and cultivate a deep connection with God in a form that resonates with our love and devotion?
\end{mananam}
\sreflect
\begin{inspiration}{\inspirationTitleFont{Inspiration}}
\normalfont\mananamtext Within the Indian tradition, various deities represent diverse forms of the One Supreme Lord. A devotee is free to choose and connect with a specific form of the Divine that is dear and pleasing. At first, this relationship may feel one-sided, as if it exists only in the heart of the devotee. Yet, through steadfast faith and devotion, the Supreme reveals Himself in the chosen form, offering solace during times of sorrow and strength in moments of weakness.
This personal relationship with the Divine stands out as the most enriching and fulfilling bond a human can experience. Across countless ages and traditions, devotees have regarded the gift to “see” God in a chosen form—one with whom they can intimately interact—as the pinnacle of spiritual achievement.
\end{inspiration}
\newpage

\slcol{\Index{नाहं वेदैर्न तपसा} न दानेन न चेज्यया । \\
शक्य एवंविधो द्रष्टुं दृष्टवानसि मां यथा ॥ ११.५३ ॥ }
\translate{nāhaṁ vēdairna tapasā na dānēna na cējyayā | \\
śakya ēvaṁvidhō draṣṭuṁ dr̥ṣṭavānasi māṁ yathā || 11.53 || }
\cquote{Neither by the study of the \emph{vedas}, nor by austerity, nor by charity and nor by sacrifices, is it possible to see Me in this form as you have seen.}
\slcol{\Index{भक्त्या त्वनन्यया} शक्य अहमेवंविधोऽर्जुन । \\
ज्ञातुं द्रष्टुं च तत्त्वेन प्रवेष्टुं च परन्तप ॥ ११.५४ ॥ }
\translate{bhaktyā tvananyayā śakya ahamēvaṁvidhō:'rjuna | 
jñātuṁ draṣṭuṁ ca tattvēna pravēṣṭuṁ ca parantapa || 11.54 || }
\cquote{But by single-minded devotion, O Arjuna, I can be seen in this way, truly known, and entered into, O Parantapa (scorcher of foes, Arjuna).}
\slcol{\Index{मत्कर्मकृन्मत्परमो} मद्भक्तः सङ्गवर्जितः । \\
निर्वैरः सर्वभूतेषु यः स मामेति पाण्डव ॥ ११.५५ ॥ }
\translate{matkarmakr̥nmatparamō madbhaktaḥ saṅgavarjitaḥ | \\
nirvairaḥ sarvabhūtēṣu yaḥ sa māmēti pāṇḍava || 11.55 || }
\cquote{He who works for Me, who regards Me as supreme, who is devoted to Me, free from attachment, and without enmity toward any being—such a one comes to Me, O Pāṇḍava (son of Pāṇḍu, Arjuna).}

\newpage
53-54-55
Mananam:
\begin{mananam}{\mananamfont {Mananam Śloka - 53, 54, 55}}
\normalfont\mananamtext Do I believe in the concept of God’s grace? If yes, why? Or, do I hold the belief that every outcome is strictly the result of one’s own actions?
Is it possible for us to witness the cosmic vision of the Supreme Lord as Arjuna did? Why is it that such a divine blessing can only be granted by God and not achieved solely through one’s past merits or virtuous deeds?
Reflecting on Arjuna as a devotee, what qualities or attitudes might have enabled him to receive this extraordinary vision? 
\end{mananam}
\sreflect
\begin{inspiration}{\inspirationTitleFont{Inspiration}}
\normalfont\mananamtext The greatest marvels of the world are limited when compared to their Creator, who can simultaneously manifest infinite forms while remaining formless in essence. In the presence of the cosmic body of the Supreme Being, our physical body is but an infinitesimal fragment. To witness the infinite through our finite senses requires a unique blessing, one that only the Lord can grant.
The Supreme Being can be truly perceived or realized solely through divine grace and consent, utilising a special set of senses far subtler than the gross physical ones. For a devotee who humbly surrenders and yearns for such a vision, the path lies in seeking God’s grace with unwavering faith and devotion.
\end{inspiration}
\newpage

\begin{center}
\devfont ॐ तत्सदिति श्रीमद्भगवद्नीतासूपनिषत्सु  ब्रह्मविद्यायां  योगशास्त्रे\\
श्रीकृष्णार्जुनसंवादे  विश्वरूप दरशन योगो नाम ऎकादशमोऽध्यायः  ॥\\
\chapEndSlokaTranslate{ōṃ tatsaditi śrīmadbhagavadnītās ūpaniṣatsu\\  brahmavidyāyāṁ  yōgaśāstrē\\
śrīkr̥ṣṇārjunasaṁvādē \\ viśvarūpa daraśana yōgō nāma ekādaśamō:'dhyāyaḥ  ||}
\end{center}

